\section{What a certificate proves}
\label{j:cert:section}

A certificate in this proof is finite data from which a small collection
of exact inequalities can be checked. It may specify coefficients of a
polynomial majorant, a tree of parameter boxes, or a schedule of bounds
which successively restricts an extremizer. The numerical statements are
about whole regions, not sampled distributions. The mathematical
reduction must first show that every still possible extremizer maps into
one of those regions.

Here is a typical example. Suppose a preceding argument has confined the
support of the same marginal to $K$ and its moment vector to $M$.
For a threshold interval $[a,b]$, a certificate specifies a function $T$
such that
\[
 T(x_1,\ldots,x_n)\ge\ind_{\{\sum_i x_i\le b\}}
 \quad(x_1,\ldots,x_n\in K),\qquad
 \E T(X_1,\ldots,X_n)\le L_a+\frac{c\rho}{\sqrt n},
\]
where the second inequality holds for every moment vector in $M$ and
$L_a\le\Phi(a/\sqrt n)$. For every $H\in[a,b]$, monotonicity gives
\[
 \Pp(S_n\le H)-\Phi(H/\sqrt n)
 \le\E T-L_a\le c\rho/\sqrt n.
\]
Thus verification proves a probability inequality, including threshold
atoms. Choosing $a$ for the normal lower bound and $b$ for the closed
event is essential. This is how the small-summand certificates work;
the count comparisons apply the same principle after conditioning on
the signs of the summands.

\subsection{A common acceptance argument}

\begin{lemma}[Expectation under one iid marginal]\label{j:cert:tensor}
Let $f_0=1,f_1,\ldots,f_d$ be absolutely integrable under $P$, and put
$m_j=\E f_j(X)$. For a finite tensor sum
\[
 T(x_1,\ldots,x_n)=\sum_\alpha t_\alpha
                         \prod_{i=1}^n f_{\alpha_i}(x_i),
 \qquad M_T(m)=\sum_\alpha t_\alpha\prod_i m_{\alpha_i},
\]
one has $\E T(X_1,\ldots,X_n)=M_T(m)$. If $T$ majorizes the
closed sum event on the support product, a uniform upper bound on
$M_T$ gives the corresponding upper bound on its probability.
\end{lemma}
\begin{proof}
Independence factors the absolute expectation of each product into a
product of finite absolute moments. Fubini and linearity therefore
evaluate every term. Integrating the pointwise majorant proves the
last assertion.
\end{proof}
The same $m$ occurs in every factor. In particular, a factor with zero
mean makes its tensor term have zero expectation. Such terms can improve
the pointwise inequality without changing its expected value. Features
of degree above three are used only after support has been bounded, so
their integrability does not strengthen the theorem's moment assumption.

\begin{lemma}[Exact positivity on a box]\label{j:cert:bernstein}
Let $p(t)=\sum_\alpha a_\alpha t^\alpha$ have coordinate degrees
at most $d_i$. Its tensor Bernstein coefficients on $[0,1]^r$ are
\[
 b_\nu=\sum_{\alpha\le\nu}a_\alpha
       \prod_{i=1}^r\frac{\binom{\nu_i}{\alpha_i}}{\binom{d_i}{\alpha_i}}.
\]
If all $b_\nu\ge0$, then $p\ge0$ on the entire closed box.
\end{lemma}
\begin{proof}
The binomial theorem gives
\[
 t^k=\sum_{j=k}^d\frac{\binom jk}{\binom dk}
                       \binom dj t^j(1-t)^{d-j}.
\]
Apply this in each coordinate. The resulting basis functions are
nonnegative and sum to one, so $p(t)$ is a convex combination of
the $b_\nu$.
\end{proof}
An affine rational substitution treats any rational box. On a simplex,
the analogous multinomial basis in barycentric coordinates is
nonnegative and sums to one, giving the same criterion after an exact
change of basis. Subdivision can improve an inconclusive coefficient
test. It never turns that test into sampling.

\begin{proposition}[Coverage and exclusion]\label{j:cert:cover}
Suppose an analytic argument maps every law and threshold under
consideration into a finite union of closed root cells. At each root,
let a finite tree subdivide a parent into closed children whose union
contains the parent. On each leaf suppose either a proved necessary
condition excludes every actual mapped point, or a valid inequality
establishes the desired score bound. Then the bound holds for every
mapped law and threshold. If all leaves exclude a putative violating
extremizer, no such extremizer exists.
\end{proposition}
\begin{proof}
A point in a root belongs to a child at each subdivision, including on
a shared face. Finiteness leads to a leaf, where its conclusion
applies. The root cover includes every actual point.
\end{proof}

There are three distinct obligations in this proposition: the map from
laws to parameters, validity of the leaf inequalities, and coverage.
A relaxed moment domain may contain unrealizable points; this only
makes verification harder. Missing an actual point invalidates the
argument. In particular, we split sign orthants and event planes
explicitly: on $\sum x_i=b$ the required inequality is $T-1\ge0$,
not just $T\ge0$. Non-strict inequalities in the domain retain endpoint
atoms. Resource limits, unresolved leaves and failed divisions cannot
be interpreted as acceptance.

\subsection{Exact evaluation and rigorous enclosures}
\label{j:cert:arithmetic}
The acceptance programs evaluate rational numbers by integer arithmetic.
Irrational constants are enclosed between rationals with proved error
bounds. For example, alternating arctangent series in Machin's identity
\[
 \frac\pi4=4\arctan(1/5)-\arctan(1/239)
\]
bound $\pi$; integer square roots then bound its required square roots.
Integrating the Taylor expansion of $e^{-t^2/2}$ gives rational bounds
for the Gaussian CDF at rational endpoints. Taylor's remainder, rather
than a printed decimal approximation, supplies the error sign and size.
The same principle bounds exponential, trigonometric and binomial
expressions used below. Each application states its finite interval or
separate tail estimate. Products and quotients use outward bounds, with
denominator positivity established on the whole cell. Multiplying a
negative expression by a lower bound requires the corresponding reversal
of inequality; the signed comparisons retain this distinction.

These statements are proved inequalities in ordinary analysis. Their
finite arithmetic evaluation is performed by the programs. Exact
arithmetic removes round-off uncertainty, but does not by itself prove
coverage, correctness of the formulas, or their applicability to a
particular marginal. The preceding lemmas and the applications below
supply those mathematical obligations. This separation is also used in
other computer-assisted inequality proofs; formal interval verification
provides a stronger implementation assurance when available
\cite{SolovyevHales2013}. The certificate format here is specific to this
argument, and is not a standard proof-assistant format.

\subsection{How the certificates are found and checked}
\label{j:cert:discovery}
For the two-summand majorants, the unknown tensor coefficients enter
linearly. The construction program solves floating-point linear
programs imposing the proposed majorant and expectation inequalities
on an initial finite set of points. It searches for violated constraints,
adds them, and solves again. Coefficients are converted to rational
numbers and passed to the exact checker. The three-summand construction
uses the same linear-programming and separation approach, with extra
zero-expectation tensor features. Success of the sampled optimization
is only a way of proposing a polynomial; the whole-domain check is the
proof. If rationalization loses a margin, the candidate must be improved
or rejected.

For the direct local comparison, the proposed certificate is a rational
quartic $Q$. Acceptance first proves that $Q$ majorizes the staircase
payoff of the coupled counts, including its jump points and both
unbounded rays. This is a univariate polynomial sign problem; exact
Sturm sequences count roots on compactified intervals. The score test
then bounds $\E Q$ using the actual conditional moments. Concavity
reduces the noise and moment variables to a finite set of corners.
The resulting polynomials in the probability and signed phase are
checked by Lemma~\ref{j:cert:bernstein} on complete closed cells.
Section~\ref{j:local:direct-finite} proves this reduction, including
the case of a negative third-moment correction.

Floating-point linear programming proposes the quartic coefficients.
A violated-constraint search improves them before rationalization;
if an exact test fails, construction may subdivide the probability or
phase interval. Acceptance reads the resulting tree and recomputes every
inequality. No sampled constraint or optimizer status establishes a leaf.
The central large-index proof instead uses an explicit quartic family,
with four exact polynomial checks printed in the text.

Resource certificates use adaptive subdivision and successive proved
enclosures. Starting with the licensed root domains, each update retains
the same law, discards an impossible cell, or splits a cell whose bound
is undecided. A successful finite forest contains a complete cover and
enough data to check every leaf. Some modules recompute the enclosure
schedule during acceptance. The schedule may have been chosen by
experimentation; soundness depends on the stated inequalities and
coverage, without requiring that the search always terminate.

The distributed candidates are treated as untrusted data. Acceptance
reconstructs the defined polynomials or bound expressions, checks their
identities and complete domains, and evaluates the exact inequalities.
It does not accept a stored success flag or an optimizer's declaration
of infeasibility. The retained public count candidates were verified in this
way; their historical search was not rerun. This has no bearing on the
logical validity of a verified certificate, just as finding a factor
and checking its product are different tasks.

\subsection{Proof data and reproduction}
\label{j:cert:reproduction}
The companion reproduction package identifies each certificate family,
its mathematical conclusion, checking source and complete execution.
The fixed reference is commit
\texttt{c31d0cbe2af867d26871fbae37a3a7f4066f4a55} of the
\href{https://github.com/kolodziejek/berry-esseen-certificates/tree/c31d0cbe2af867d26871fbae37a3a7f4066f4a55/manuscript/extended-v17/reproduction}
{public certificate repository}, in
\texttt{manuscript/extended-v17/reproduction}.
Its \texttt{REPRODUCE.md} gives the restoration and acceptance commands;
\texttt{COMPLETION\_LEDGER.json} identifies the proof objects and
completed result records. The mathematical interfaces used in this
revision are stated in the numerical propositions below.
This reference preserves the original acquisition domains and the
unchanged finite-base and Fourier families. The direct-comparison
supplement contains the new quartics, complete domain catalogues and
checking programs. Its reproduction instructions identify the selected
source and proof-object hashes. The finite comparisons retain each
original allowance, actual index and weighted resource.
The supplement is distributed with this revision in
\texttt{certificates/direct-comparison}. Its three archives preserve the
complete objects and source inventories for the earlier, left and central
comparisons; their README files give exact replay commands. The central
archive also includes the explicit uniform-tail calculation.
A reader can inspect the analytic arguments independently of running
the programs, and can verify the finite arithmetic independently of
repeating the search. Both parts are necessary for the whole proof.
The completed executions use Python with assertions enabled, with integer and rational
arithmetic; the three-summand kernel uses C++17 and
Boost.Multiprecision integers. The compiler, language runtimes and
faithful execution are part of the trusted computational environment.
The implementation is not a proof-assistant formalization.

The full execution history, command lines and machine identities belong
to that package rather than the mathematical argument. They distinguish
complete native calculations from later checks of their input
interfaces, and preserve failed or interrupted attempts. The selected numerical obligations of the reference proof retain their
completed records; the separate 114-command recipe has not been rerun as
one relocated campaign. A fresh replay of a new object uses its exact
acceptance program and is recorded separately from construction. Replaying
the same program does not constitute an independent implementation.
The new direct comparisons replace their selected predecessors; the
other certificate bytes and targets are retained. Any extension of a
certificate's scope must be justified explicitly and cannot follow from
its previous success alone.
